%==============================================================================%
% Chapter 7 : GPU (CUDA) computation
%==============================================================================%
\chapter{\bi{GPU computation with CUDA}{CUDAによるGPU計算}}\label{chap:gpu}

\bi{This chapter describes the CUDA implementation of BNCmatmul, which executes the
basic linear computation, matrix multiplication, LU decomposition and sparse
matrix-vector multiplication on an NVIDIA GPU.  It is built on top of
\textbf{gdtq-0.0.2} (the device multi-component floating-point header layer) and,
for arbitrary precision, on \textbf{mpc\_cuda-0.0.1} (a device port of MPFR/MPC).
The reference target is an NVIDIA GB10 (\texttt{sm\_121}, CUDA 13).}{本章では，
基本的な線形計算，行列乗算，LU分解，および疎行列ベクトル乗算をNVIDIA GPU上で実行する
BNCmatmulのCUDA実装について説明します。本実装は\textbf{gdtq-0.0.2}（デバイス用多成分
浮動小数点ヘッダ層）の上に構築されており，任意精度については\textbf{mpc\_cuda-0.0.1}
（MPFR/MPCのデバイス移植版）を用います。基準とする対象はNVIDIA GB10
（\texttt{sm\_121}，CUDA 13）です。}

\bi{All GPU functions are gathered into the static library
\texttt{libbncmatmul-0.24\_cuda.a}.  The single script
\texttt{bench/build\_cuda\_full.sh} compiles every GPU module into this library and
runs the whole verification suite (matrix multiplication, matrix-vector product,
LU and SpMV for every precision, plus the CUDA MPFR/MPC matrix multiplication).}{すべての
GPU関数は静的ライブラリ\texttt{libbncmatmul-0.24\_cuda.a}にまとめられています。単一の
スクリプト\texttt{bench/build\_cuda\_full.sh}は，すべてのGPUモジュールをこのライブラリに
コンパイルし，検証スイート全体（各精度の行列乗算，行列ベクトル積，LU，SpMV，および
CUDA MPFR/MPC行列乗算）を実行します。}

%------------------------------------------------------------------%
\section{\bi{Naming convention and data types}{命名規則とデータ型}}
%------------------------------------------------------------------%

\bi{Every GPU type and function carries a leading \textsf{g} (``GPU'').  The element
type follows the same multi-component scheme as the CPU side: a \textsf{d}-family
based on \texttt{double} and an \textsf{s}-family based on \texttt{float}.}{すべての
GPU型および関数には先頭に\textsf{g}（``GPU''）が付きます。要素型はCPU側と同じ多成分方式に
従い，\texttt{double}に基づく\textsf{d}系列と，\texttt{float}に基づく\textsf{s}系列が
あります。}

\begin{table}[h]
\centering
\begin{tabular}{l l l l}
\hline
GPU prefix & element type & bits ($\approx$) & CPU counterpart \\
\hline
\textsf{gf}  & \texttt{float} (\texttt{float})            & 24  & FVector / FMatrix \\
\textsf{gd}  & \texttt{double} (\texttt{double})          & 53  & DVector / DMatrix \\
\textsf{gds} & double-single (\texttt{float2})            & 48  & DSVector / DSMatrix \\
\textsf{gts} & triple-single (\texttt{float3})            & 72  & TSVector / TSMatrix \\
\textsf{gqs} & quadruple-single (\texttt{float4})         & 96  & QSVector / QSMatrix \\
\textsf{gdd} & double-double (\texttt{double2})           & 106 & DDVector / DDMatrix \\
\textsf{gtd} & triple-double (\texttt{double3})           & 159 & TDVector / TDMatrix \\
\textsf{gqd} & quadruple-double (\texttt{double4})        & 212 & QDVector / QDMatrix \\
\hline
\end{tabular}
\caption{\bi{Real-valued GPU element types.  \textsf{g$\tau$} with
$\tau\in\{$f, d, ds, ts, qs, dd, td, qd$\}$.}{実数値のGPU要素型。\textsf{g$\tau$}は
$\tau\in\{$f, d, ds, ts, qs, dd, td, qd$\}$とする。}}
\end{table}

\bi{For complex numbers two families are provided.  The native complex types
\textsf{gcd} (complex \texttt{double}) and \textsf{gcf} (complex \texttt{float})
store the real and imaginary parts as two separate scalar arrays.  The
multi-component complex types \textsf{cgdd}, \textsf{cgtd}, \textsf{cgqd},
\textsf{cgds}, \textsf{cgts}, \textsf{cgqs} store the real and imaginary parts as
two separate arrays of the corresponding \textsf{g$\tau$} element.}{複素数には2つの
系列が用意されています。ネイティブ複素型\textsf{gcd}（複素\texttt{double}）および
\textsf{gcf}（複素\texttt{float}）は，実部と虚部を2つの別々のスカラ配列として格納します。
多成分複素型\textsf{cgdd}，\textsf{cgtd}，\textsf{cgqd}，\textsf{cgds}，\textsf{cgts}，
\textsf{cgqs}は，実部と虚部を対応する\textsf{g$\tau$}要素の2つの別々の配列として格納します。}

\bi{Each family defines a vector and a matrix object.  For a real type
\textsf{g$\tau$}:}{各系列はベクトルオブジェクトと行列オブジェクトを定義します。実数型
\textsf{g$\tau$}の場合は次のとおりです。}
\begin{verbatim}
typedef struct { long int dim;                gXX_real *element; } gXXvector;
typedef struct { long int row_dim, col_dim;   gXX_real *element; } gXXmatrix;
\end{verbatim}
\bi{and for the complex types the single \texttt{element} pointer is replaced by a
pair \texttt{re}, \texttt{im}.  A pointer to the structure is named
\textsf{GXXVector} / \textsf{GXXMatrix} (upper case), e.g.\ \textsf{GDDVector},
\textsf{GCDMatrix}, \textsf{CGTSVector}.}{複素型の場合は，単一の\texttt{element}
ポインタが\texttt{re}，\texttt{im}の対に置き換えられます。構造体へのポインタは
\textsf{GXXVector} / \textsf{GXXMatrix}（大文字）と名付けられます。例えば
\textsf{GDDVector}，\textsf{GCDMatrix}，\textsf{CGTSVector}などです。}

\bi{Matrices are stored row-major with stride equal to \texttt{col\_dim} (the GPU side
does \emph{not} use the SIMD column padding of the CPU side).  Two block/grid
arguments \varname{int}{num\_blocks\_per\_grid} and
\varname{int}{num\_threads\_per\_block} are passed to every kernel-launching
function; typical values are $32$--$128$ each.}{行列は行優先で，ストライドは
\texttt{col\_dim}に等しい形で格納されます（GPU側はCPU側のSIMD列パディングを
\emph{使用しません}）。2つのブロック／グリッド引数\varname{int}{num\_blocks\_per\_grid}
および\varname{int}{num\_threads\_per\_block}が，カーネルを起動するすべての関数に
渡されます。典型的な値はそれぞれ$32$～$128$です。}

%------------------------------------------------------------------%
\section{\bi{Device memory and host-device transfer}{デバイスメモリとホスト・デバイス間転送}}
%------------------------------------------------------------------%

\bi{In the list below \textsf{g$\tau$} stands for any of the real GPU types and
\textsf{XX} for its CPU counterpart (e.g.\ \textsf{gdd}$\leftrightarrow$\textsf{DD}).}{以下の
リストにおいて，\textsf{g$\tau$}は実数GPU型のいずれかを表し，\textsf{XX}はそのCPU側の
対応物を表します（例：\textsf{gdd}$\leftrightarrow$\textsf{DD}）。}

\begin{itemize}
\funcitem{GXXVector}{init\_gXXvector}{\varnameli{dim}}\ $\cdots$ \bi{Allocate a vector in
host memory.}{ホストメモリにベクトルを確保します。}
\funcitem{GXXVector}{init\_gXXvector\_dev}{\varnameli{dim}}\ $\cdots$ \bi{Allocate a vector
in device (GPU) memory.}{デバイス（GPU）メモリにベクトルを確保します。}
\funcitem{void}{free\_gXXvector}{\varname{GXXVector}{vec}}
\funcitem{void}{free\_gXXvector\_dev}{\varname{GXXVector}{vec}}
\funcitem{GXXMatrix}{init\_gXXmatrix}{\varnameli{row\_dim}, \varnameli{col\_dim}}
\funcitem{GXXMatrix}{init\_gXXmatrix\_dev}{\varnameli{row\_dim}, \varnameli{col\_dim}}
\funcitem{void}{free\_gXXmatrix\_dev}{\varname{GXXMatrix}{mat}}
\funcitem{void}{subst\_gXXvector\_dev\_XXvec}{\varname{GXXVector}{dev}, \varname{XXVector}{src}}\
$\cdots$ \bi{Copy a CPU vector to the GPU (host $\to$ device).}{CPUのベクトルをGPUへ
コピーします（ホスト$\to$デバイス）。}
\funcitem{void}{subst\_XXvector\_gXXvec\_dev}{\varname{XXVector}{dst}, \varname{GXXVector}{dev}}\
$\cdots$ \bi{Copy a GPU vector back to the CPU (device $\to$ host).}{GPUのベクトルを
CPUへ戻します（デバイス$\to$ホスト）。}
\funcitem{void}{subst\_gXXmatrix\_dev\_XXmat}{\varname{GXXMatrix}{dev}, \varname{XXMatrix}{src}}
\funcitem{void}{subst\_XXmatrix\_gXXmat\_dev}{\varname{XXMatrix}{dst}, \varname{GXXMatrix}{dev}}
\funcitem{void}{print\_gXXvector\_dev}{\varname{GXXVector}{dev}}
\funcitem{void}{print\_gXXmatrix\_dev}{\varname{GXXMatrix}{dev}}
\end{itemize}

\bi{For the native complex types \textsf{gcd}/\textsf{gcf} the CPU counterpart is the
complex-\texttt{double} \textsf{CDVector}/\textsf{CDMatrix}; the transfer
functions are \funcname{subst\_gcdvector\_dev\_cdvec}{} and
\funcname{subst\_cdvector\_gcdvec\_dev}{} (and the matrix variants).  For the
multi-component complex types \textsf{cg$\tau$} the CPU counterpart is the
complex multi-component \textsf{CXXVector}/\textsf{CXXMatrix} (e.g.\
\textsf{CDDMatrix}), with \funcname{subst\_cgddvector\_dev\_cddvec}{} etc.}{ネイティブ
複素型\textsf{gcd}/\textsf{gcf}のCPU側対応物は，複素\texttt{double}の
\textsf{CDVector}/\textsf{CDMatrix}です。転送関数は
\funcname{subst\_gcdvector\_dev\_cdvec}{}および
\funcname{subst\_cdvector\_gcdvec\_dev}{}（および行列版）です。多成分複素型
\textsf{cg$\tau$}のCPU側対応物は，複素多成分の\textsf{CXXVector}/\textsf{CXXMatrix}
（例：\textsf{CDDMatrix}）であり，\funcname{subst\_cgddvector\_dev\_cddvec}{}などを
用います。}

%------------------------------------------------------------------%
\section{\bi{Vector and matrix arithmetic}{ベクトルと行列の演算}}
%------------------------------------------------------------------%

\bi{All routines below run on the GPU and take the two launch arguments
\varname{int}{nbg} (blocks) and \varname{int}{ntb} (threads); they are omitted
from the argument lists for brevity.}{以下のルーチンはすべてGPU上で動作し，2つの起動引数
\varname{int}{nbg}（ブロック数）と\varname{int}{ntb}（スレッド数）を取ります。
簡潔さのため，これらは引数リストから省略しています。}

\begin{itemize}
\funcitem{void}{add\_gXXvector\_dev}{\varname{GXXVector}{c}, \varname{GXXVector}{a}, \varname{GXXVector}{b}}\
$\cdots$ $\mathbf{c} := \mathbf{a}+\mathbf{b}$.
\funcitem{void}{sub\_gXXvector\_dev}{\varname{GXXVector}{c}, \varname{GXXVector}{a}, \varname{GXXVector}{b}}\
$\cdots$ $\mathbf{c} := \mathbf{a}-\mathbf{b}$.
\funcitem{void}{cmul\_gXXvector\_dev}{\varname{GXXVector}{c}, \varname{gXX\_real}{val}, \varname{GXXVector}{a}}\
$\cdots$ $\mathbf{c} := \mathrm{val}\cdot\mathbf{a}$.
\funcitem{void}{subst\_gXXvector\_dev}{\varname{GXXVector}{c}, \varname{GXXVector}{a}}\ $\cdots$ $\mathbf{c} := \mathbf{a}$.
\funcitem{void}{set0\_gXXvector\_dev}{\varname{GXXVector}{c}}\ $\cdots$ $\mathbf{c} := \mathbf{0}$.
\funcitem{void}{ip\_gXXvector\_dev}{\varname{gXX\_real*}{ret\_dev}, \varname{GXXVector}{a}, \varname{GXXVector}{b}}\
$\cdots$ \bi{inner product $(\mathbf{a},\mathbf{b})$, result written to device scalar \texttt{ret\_dev}.}{内積$(\mathbf{a},\mathbf{b})$。結果はデバイススカラ\texttt{ret\_dev}に書き込まれます。}
\funcitem{void}{norm2\_gXXvector\_dev}{\varname{gXX\_real*}{ret\_dev}, \varname{GXXVector}{a}}\ $\cdots$ $\|\mathbf{a}\|_2$.
\funcitem{void}{norm1\_gXXvector\_dev}{\varname{gXX\_real*}{ret\_dev}, \varname{GXXVector}{a}}\ $\cdots$ $\|\mathbf{a}\|_1$.
\funcitem{void}{normi\_gXXvector\_dev}{\varname{gXX\_real*}{ret\_dev}, \varname{GXXVector}{a}}\ $\cdots$ $\|\mathbf{a}\|_\infty$.

\funcitem{void}{add\_gXXmatrix\_dev}{\varname{GXXMatrix}{C}, \varname{GXXMatrix}{A}, \varname{GXXMatrix}{B}}\ $\cdots$ $C := A+B$.
\funcitem{void}{sub\_gXXmatrix\_dev}{\varname{GXXMatrix}{C}, \varname{GXXMatrix}{A}, \varname{GXXMatrix}{B}}\ $\cdots$ $C := A-B$.
\funcitem{void}{cmul\_gXXmatrix\_dev}{\varname{GXXMatrix}{C}, \varname{gXX\_real}{sc}, \varname{GXXMatrix}{A}}\ $\cdots$ $C := \mathrm{sc}\cdot A$.
\funcitem{void}{transpose\_gXXmatrix\_dev}{\varname{GXXMatrix}{C}, \varname{GXXMatrix}{A}}\ $\cdots$ $C := A^{T}$.
\funcitem{void}{subst\_gXXmatrix\_dev}{\varname{GXXMatrix}{C}, \varname{GXXMatrix}{A}}\ $\cdots$ $C := A$.
\funcitem{void}{set0\_gXXmatrix\_dev}{\varname{GXXMatrix}{C}}\ $\cdots$ $C := O$.
\funcitem{void}{setI\_gXXmatrix\_dev}{\varname{GXXMatrix}{C}}\ $\cdots$ $C := I$.
\funcitem{void}{normf\_gXXmatrix\_dev}{\varname{gXX\_real*}{ret\_dev}, \varname{GXXMatrix}{A}}\ $\cdots$ \bi{Frobenius norm $\|A\|_F$.}{フロベニウスノルム$\|A\|_F$。}
\end{itemize}

\bi{The complex families additionally provide
\funcname{conjtrans\_gcdmatrix\_dev}{} / \funcname{conjtrans\_cgXXmatrix\_dev}{}
($C := A^{H}$, conjugate transpose), and the complex scalar multiplication takes
the real and imaginary parts as two arguments, e.g.\
\funcname{cmul\_gcdvector\_dev}{\varname{GCDVector}{c}, \varname{double}{val\_re}, \varname{double}{val\_im}, \varname{GCDVector}{a}}.}{複素系列は
さらに\funcname{conjtrans\_gcdmatrix\_dev}{} / \funcname{conjtrans\_cgXXmatrix\_dev}{}
（$C := A^{H}$，共役転置）を提供します。また複素スカラ乗算は実部と虚部を2つの引数として
取ります。例：\funcname{cmul\_gcdvector\_dev}{\varname{GCDVector}{c}, \varname{double}{val\_re}, \varname{double}{val\_im}, \varname{GCDVector}{a}}。}

%------------------------------------------------------------------%
\section{\bi{Matrix multiplication and matrix-vector product}{行列乗算と行列ベクトル積}}
%------------------------------------------------------------------%

\begin{itemize}
\funcitem{void}{mul\_gXXmatrix\_dev}{\varname{GXXMatrix}{C}, \varname{GXXMatrix}{A}, \varname{GXXMatrix}{B}}\
$\cdots$ \bi{$C := AB$ (one thread per output row, grid-stride).}{$C := AB$（出力行ごとに1スレッド，グリッドストライド）。}
\funcitem{void}{mul\_gXXmatrix\_gXXvec}{\varname{GXXVector}{v}, \varname{GXXMatrix}{A}, \varname{GXXVector}{vb}}\
$\cdots$ $\mathbf{v} := A\mathbf{vb}$.
\funcitem{void}{mul\_gXXmatrixt\_gXXvec}{\varname{GXXVector}{v}, \varname{GXXMatrix}{A}, \varname{GXXVector}{vb}}\
$\cdots$ $\mathbf{v} := A^{T}\mathbf{vb}$.
\end{itemize}

\bi{For the complex families the same three routines exist, e.g.\
\funcname{mul\_gcdmatrix\_dev}{}, \funcname{mul\_cgddmatrix\_dev}{},
\funcname{mul\_cgddmatrix\_cgddvec}{} and \funcname{mul\_cgddmatrixt\_cgddvec}{};
complex multiplication uses the $4$-real-multiply formula
$\mathrm{Re}(AB)=A_rB_r-A_iB_i$, $\mathrm{Im}(AB)=A_rB_i+A_iB_r$.}{複素系列にも同じ3つの
ルーチンが存在します。例：\funcname{mul\_gcdmatrix\_dev}{}，
\funcname{mul\_cgddmatrix\_dev}{}，\funcname{mul\_cgddmatrix\_cgddvec}{}，
\funcname{mul\_cgddmatrixt\_cgddvec}{}。複素乗算は4回の実数乗算による公式
$\mathrm{Re}(AB)=A_rB_r-A_iB_i$，$\mathrm{Im}(AB)=A_rB_i+A_iB_r$を用います。}

\bi{A Strassen matrix multiplication for the multi-component types is provided in
\texttt{src/matmul\_strassen\_general\_g\{dd,td,qd,ds,ts,qs\}.cu}.}{多成分型に対する
Strassen行列乗算が
\texttt{src/matmul\_strassen\_general\_g\{dd,td,qd,ds,ts,qs\}.cu}に提供されています。}

\bi{Verified accuracy of $C:=AB$ on the GB10 against the corresponding CPU routine:
\textsf{gf} and the quadruple-single \textsf{cgqs} are bit-identical
($0$); \textsf{gd} $\sim10^{-16}$, \textsf{gcd} $\sim10^{-19}$, \textsf{gdd}
$\sim10^{-31}$, \textsf{gds} $\sim10^{-14}$, \textsf{gts} $\sim10^{-22}$.}{GB10上で
対応するCPUルーチンと比較して検証した$C:=AB$の精度は次のとおりです。\textsf{gf}および
4倍単精度の\textsf{cgqs}はビット単位で一致（$0$）。\textsf{gd}は$\sim10^{-16}$，
\textsf{gcd}は$\sim10^{-19}$，\textsf{gdd}は$\sim10^{-31}$，\textsf{gds}は$\sim10^{-14}$，
\textsf{gts}は$\sim10^{-22}$です。}

%------------------------------------------------------------------%
\section{\bi{LU decomposition and linear solver on GPU}{GPU上のLU分解と線形ソルバ}}
%------------------------------------------------------------------%

\bi{A right-looking LU decomposition with partial pivoting and a forward/back
substitution solver are provided for every real type and for the native complex
types.  The decomposition and the solve are combined in one driver:}{部分ピボット選択を
伴う右向き（right-looking）LU分解と，前進／後退代入ソルバが，すべての実数型および
ネイティブ複素型に対して提供されています。分解と求解は1つのドライバにまとめられています。}

\begin{itemize}
\funcitem{int}{gXX\_lu\_solve\_dev}{\varname{GXXMatrix}{A}, \varname{GXXVector}{b}, \varname{GXXVector}{x}, \varnameli{ch[]}, \varname{int}{blocks}, \varname{int}{threads}}\
$\cdots$ \bi{Overwrite $A$ with its $LU$ factors and solve $A\mathbf{x}=\mathbf{b}$
into $\mathbf{x}$ (all device objects).  The pivot row chosen at each step is
recorded in \texttt{ch[]} (host, length $n$; may be \texttt{NULL}).  Returns $0$ on
success.}{$A$をその$LU$因子で上書きし，$A\mathbf{x}=\mathbf{b}$を解いて
$\mathbf{x}$に格納します（すべてデバイスオブジェクト）。各ステップで選択された
ピボット行は\texttt{ch[]}（ホスト，長さ$n$。\texttt{NULL}でもよい）に記録されます。
成功時に$0$を返します。}
\end{itemize}

\bi{\texttt{XX} is one of \texttt{d}, \texttt{f}, \texttt{dd}, \texttt{td}, \texttt{qd},
\texttt{ds}, \texttt{ts}, \texttt{qs} (real, e.g.\ \funcname{gqs\_lu\_solve\_dev}{}),
or \texttt{cd}, \texttt{cf} (native complex, \funcname{gcd\_lu\_solve\_dev}{} /
\funcname{gcf\_lu\_solve\_dev}{}).  All ten drivers are declared
\texttt{extern "C"}.  Verified maximum relative error against the known solution
(dimension $100$): \textsf{gd} $1.2\times10^{-12}$, \textsf{gf}
$9.5\times10^{-7}$, \textsf{gcd} $2.8\times10^{-15}$, \textsf{gdd}/\textsf{gtd}/\textsf{gds} $0$,
\textsf{gts} $7.7\times10^{-22}$, \textsf{gqs} $2.8\times10^{-29}$.}{\texttt{XX}は
\texttt{d}，\texttt{f}，\texttt{dd}，\texttt{td}，\texttt{qd}，\texttt{ds}，\texttt{ts}，
\texttt{qs}（実数。例：\funcname{gqs\_lu\_solve\_dev}{}），または\texttt{cd}，\texttt{cf}
（ネイティブ複素。\funcname{gcd\_lu\_solve\_dev}{} / \funcname{gcf\_lu\_solve\_dev}{}）の
いずれかです。10個のドライバはすべて\texttt{extern "C"}として宣言されています。既知の解に
対して検証した最大相対誤差（次元$100$）は次のとおりです。\textsf{gd}は
$1.2\times10^{-12}$，\textsf{gf}は$9.5\times10^{-7}$，\textsf{gcd}は$2.8\times10^{-15}$，
\textsf{gdd}/\textsf{gtd}/\textsf{gds}は$0$，\textsf{gts}は$7.7\times10^{-22}$，
\textsf{gqs}は$2.8\times10^{-29}$です。}

%------------------------------------------------------------------%
\section{\bi{Sparse matrix-vector multiplication (SpMV)}{疎行列ベクトル乗算（SpMV）}}
%------------------------------------------------------------------%

\bi{The GPU sparse matrix is stored in the standard CSR (Compressed Sparse Row)
format, independent of the CPU ``DRS'' format.  The structure holds the
dimensions, the number of non-zeros, and three device arrays:}{GPUの疎行列は，CPUの
``DRS''形式とは独立に，標準的なCSR（Compressed Sparse Row）形式で格納されます。構造体は
次元，非ゼロ要素数，および3つのデバイス配列を保持します。}

\begin{verbatim}
typedef struct {
    long int row_dim, col_dim, nnz;
    long int *row_ptr;   /* length row_dim + 1 */
    long int *col_idx;   /* length nnz         */
    gXX_real *val;       /* length nnz         */
} gXXspmatrix;           /* GXXSPMatrix = gXXspmatrix *  */
\end{verbatim}

\begin{itemize}
\funcitem{GXXSPMatrix}{init\_gXXspmatrix\_dev}{\varnameli{row\_dim}, \varnameli{col\_dim}, \varnameli{nnz}, \varname{const long int*}{row\_ptr}, \varname{const long int*}{col\_idx}, \varname{const gXX\_real*}{val}}\
$\cdots$ \bi{Build a device CSR matrix from host CSR arrays.}{ホストのCSR配列から
デバイスのCSR行列を構築します。}
\funcitem{void}{free\_gXXspmatrix\_dev}{\varname{GXXSPMatrix}{mat}}
\funcitem{void}{mul\_gXXspmatrix\_gXXvec}{\varname{GXXVector}{ret}, \varname{GXXSPMatrix}{A}, \varname{GXXVector}{x}, \varname{int}{nbg}, \varname{int}{ntb}}\
$\cdots$ \bi{$\mathbf{ret} := A\mathbf{x}$ (one thread per row).}{$\mathbf{ret} := A\mathbf{x}$（行ごとに1スレッド）。}
\end{itemize}

\bi{The native \textsf{gd}/\textsf{gf} SpMV uses the \textsf{GDVector}/\textsf{GFVector}
dense vectors as above; the multi-component real types (\textsf{gdd}\dots
\textsf{gqs}) take raw device \texttt{gXX\_real *} arrays for $\mathbf{x}$ and
$\mathbf{ret}$ via \funcname{mul\_gXXspmatrix}{\varname{gXX\_real*}{y}, \varname{GXXSPMatrix}{A}, \varname{const gXX\_real*}{x}, \dots}.
The native complex \textsf{gcd}/\textsf{gcf} and multi-component complex
\textsf{cg$\tau$} SpMV keep the real/imaginary value arrays separate, e.g.\
\funcname{mul\_gcdspmatrix}{\varname{double*}{yre}, \varname{double*}{yim}, \varname{GCDSPMatrix}{A}, \varname{const double*}{xre}, \varname{const double*}{xim}, \dots}.
The transpose product $A^{T}\mathbf{x}$ is obtained either by the native
\funcname{mul\_gdspmatrixt\_gdvec}{} (atomic scatter) or, for the multi-component
types, by building the CSR of $A^{T}$ on the host and reusing the $A\mathbf{x}$
kernel.}{ネイティブの\textsf{gd}/\textsf{gf} SpMVは，上記のとおり
\textsf{GDVector}/\textsf{GFVector}の密ベクトルを使用します。多成分実数型
（\textsf{gdd}\dots\textsf{gqs}）は，$\mathbf{x}$と$\mathbf{ret}$に対して生のデバイス
\texttt{gXX\_real *}配列を取り，\funcname{mul\_gXXspmatrix}{\varname{gXX\_real*}{y}, \varname{GXXSPMatrix}{A}, \varname{const gXX\_real*}{x}, \dots}を介して扱います。
ネイティブ複素\textsf{gcd}/\textsf{gcf}および多成分複素\textsf{cg$\tau$}のSpMVは，
実部／虚部の値配列を分離して保持します。例：\funcname{mul\_gcdspmatrix}{\varname{double*}{yre}, \varname{double*}{yim}, \varname{GCDSPMatrix}{A}, \varname{const double*}{xre}, \varname{const double*}{xim}, \dots}。
転置積$A^{T}\mathbf{x}$は，ネイティブの\funcname{mul\_gdspmatrixt\_gdvec}{}
（アトミックスキャッタ）によって，または多成分型の場合は$A^{T}$のCSRをホスト上で構築して
$A\mathbf{x}$カーネルを再利用することによって得られます。}

%------------------------------------------------------------------%
\section{\bi{Arbitrary precision matrix multiplication (CUDA MPFR / MPC)}{任意精度行列乗算（CUDA MPFR / MPC）}}
%------------------------------------------------------------------%

\bi{Using the device port \textsf{mpc\_cuda}, the dot products of a matrix
multiplication or a matrix-vector product can be accumulated at an arbitrary
compile-time precision (default $1024$ bits) on the GPU.  The matrices and
vectors are passed as plain row-major \texttt{double} arrays (the complex
variants use separate real/imaginary \texttt{double} arrays).}{デバイス移植版
\textsf{mpc\_cuda}を用いると，行列乗算または行列ベクトル積の内積を，コンパイル時に指定した
任意精度（既定$1024$ビット）でGPU上に累算できます。行列とベクトルは，単純な行優先の
\texttt{double}配列として渡されます（複素版は実部／虚部を分離した\texttt{double}配列を
使用します）。}

\begin{itemize}
\funcitem{int}{cuda\_mul\_mpfmatrix}{\varname{double*}{C}, \varname{const double*}{A}, \varname{const double*}{B}, \varname{int}{n}, \varname{int}{lblocks}, \varname{int}{lthreads}}\
$\cdots$ \bi{$C := AB$ ($n\times n$), each entry accumulated in MPFR at \texttt{PREC} bits.}{$C := AB$（$n\times n$）。各要素を\texttt{PREC}ビットのMPFRで累算します。}
\funcitem{int}{cuda\_mul\_mpfmatrix\_vec}{\varname{double*}{y}, \varname{const double*}{A}, \varname{const double*}{x}, \varname{int}{n}, \varname{int}{lblocks}, \varname{int}{lthreads}}\
$\cdots$ $\mathbf{y} := A\mathbf{x}$.
\funcitem{int}{cuda\_mul\_cmpfmatrix}{\varname{double*}{Cr}, \varname{double*}{Ci}, \varname{const double*}{Ar}, \varname{const double*}{Ai}, \varname{const double*}{Br}, \varname{const double*}{Bi}, \varname{int}{n}, \varname{int}{lblocks}, \varname{int}{lthreads}}\
$\cdots$ \bi{Complex $C := AB$ accumulated in MPC.}{複素$C := AB$。MPCで累算します。}
\funcitem{int}{cuda\_mul\_cmpfmatrix\_vec}{\varname{double*}{yr}, \varname{double*}{yi}, \varname{const double*}{Ar}, \varname{const double*}{Ai}, \varname{const double*}{xr}, \varname{const double*}{xi}, \varname{int}{n}, \varname{int}{lblocks}, \varname{int}{lthreads}}
\end{itemize}

\bi{These routines link the static library \texttt{/usr/local/lib/libmpc\_cuda.a} and
require the device stack/heap limits to be enlarged (the drivers set
\texttt{cudaLimitStackSize} and the device heap internally).  Verified on the
GB10: both the MPFR ($n=96$) and MPC ($n=80$) matrix multiplication and
matrix-vector product agree with the CPU result to $0$ at $1024$ bits.}{これらの
ルーチンは静的ライブラリ\texttt{/usr/local/lib/libmpc\_cuda.a}をリンクし，デバイスの
スタック／ヒープ上限を拡大する必要があります（ドライバは内部で
\texttt{cudaLimitStackSize}とデバイスヒープを設定します）。GB10上で検証した結果，
MPFR（$n=96$）とMPC（$n=80$）の行列乗算および行列ベクトル積は，いずれも$1024$ビットで
CPUの結果と$0$まで一致しました。}

%------------------------------------------------------------------%
\section{\bi{Building and running the GPU suite}{GPUスイートのビルドと実行}}
%------------------------------------------------------------------%

\begin{verbatim}
sh bench/build_cuda_full.sh
\end{verbatim}
\bi{compiles every module listed above into \texttt{libbncmatmul-0.24\_cuda.a} and runs
the complete verification.  The device-math operators are supplied by the
\textsf{gdtq} aggregators: \texttt{gqd.cu} (\texttt{double2/3/4}) for the
\textsf{d}-family and \texttt{gqs.cu} (\texttt{float2/3/4}) for the
\textsf{s}-family; the \textsf{gdd}/\textsf{gds} objects already carry their
operators inline, while \textsf{gtd}/\textsf{gqd}/\textsf{gts}/\textsf{gqs} are
device-linked against the corresponding aggregator object.  The arbitrary
precision modules need the \textsf{mpc\_cuda} headers under
\texttt{/usr/local/include/mpc\_cuda} and the definitions in
\texttt{/usr/local/share/mpc\_cuda/mpfr\_cuda\_defs.txt}.}{は，上記でリストした
すべてのモジュールを\texttt{libbncmatmul-0.24\_cuda.a}にコンパイルし，完全な検証を
実行します。デバイス算術演算子は\textsf{gdtq}の集約モジュールによって供給されます。
すなわち\textsf{d}系列には\texttt{gqd.cu}（\texttt{double2/3/4}），\textsf{s}系列には
\texttt{gqs.cu}（\texttt{float2/3/4}）です。\textsf{gdd}/\textsf{gds}オブジェクトは
すでに演算子をインラインで保持していますが，\textsf{gtd}/\textsf{gqd}/\textsf{gts}/\textsf{gqs}は
対応する集約オブジェクトに対してデバイスリンクされます。任意精度モジュールは，
\texttt{/usr/local/include/mpc\_cuda}以下の\textsf{mpc\_cuda}ヘッダと，
\texttt{/usr/local/share/mpc\_cuda/mpfr\_cuda\_defs.txt}内の定義を必要とします。}
